\chapter{Non-Linear Normal Form}\label{chap:nonlinear-normal-form}

The linear normal form separates the geometry of an oscillation from its
phase advance: the matrix $A$ maps circles into phase-space ellipses, while
$R$ rotates the normalised coordinates.  The nonlinear construction makes
the same separation using canonical maps.  A nonlinear change of coordinates
describes the deformation of the invariant tori, and the transformed
one-turn map describes the dynamics on them.  Away from resonances, this
dynamics is an action-dependent rotation.  Near a resonance, some
angle-dependent terms must remain and can exchange action between modes.

This chapter develops the construction locally about an elliptic closed
orbit.  Nonlinear maps and their inverses are understood as formal Taylor
series, or as series truncated at a stated order.  An action-only normal
form at finite order is an approximation to the original map; it does not
assert that an arbitrary nonlinear lattice is globally integrable.  The
map-based viewpoint is also the basis of the generalised Courant--Snyder
theory described in Ref.~\cite{Forest_book2}.



\section{The next step after linear normalisation}

\subsection{Centre and normalise before expanding}

Let $\mathcal M$ be the one-turn point map at a fixed observation location,
and let $\vec z_{\mathrm{co}}$ be its closed orbit.  Use the definitions
\begin{equation}
    \mathcal T_{\mathrm{co}}(\vec\xi) = \vec z_{\mathrm{co}} + \vec\xi,
    \qquad \mathcal L_A(\vec w) = A\vec w\,,
\end{equation}
where $A$ is the symplectic linear normalising matrix of the previous
chapter.  The map in linearly normalised coordinates is
\begin{equation}\label{eq:nnf-linear-normalisation}
    \mathcal F = \mathcal L_A^{-1}\circ\mathcal T_{\mathrm{co}}^{-1}
    \circ\mathcal M\circ\mathcal T_{\mathrm{co}}\circ\mathcal L_A,
    \qquad
    \mathcal F(\vec w) = R\vec w + \mathcal O(\|\vec w\|^2)\,.
\end{equation}
Thus $\mathcal F(0) = 0$ and $D\mathcal F(0) = R$.  In particular, the
nonlinear terms describe deviations from the \emph{actual closed orbit},
not deviations from the design trajectory.  Expanding about the wrong
origin leaves a constant term and gives the wrong linear optics through
feed-down.  Recentring is equally necessary for every momentum offset or
magnet setting considered below.

We seek a canonical near-identity map
\begin{equation}
    \vec w = \mathcal B(\vec u),\qquad
    \mathcal B(0) = 0,\qquad D\mathcal B(0) = I\,,
\end{equation}
Here $\mathcal B$ is a nonlinear normalising map.  The matrix $A$ has
already removed the linear optical geometry, while $\mathcal B$ removes
nonresonant nonlinear harmonics order by order.

We require
\begin{equation}\label{eq:nnf-conjugacy}
    \mathcal N = \mathcal B^{-1}\circ\mathcal F\circ\mathcal B\,.
\end{equation}
The complete transformation from nonlinear normal coordinates to physical
coordinates is therefore
\begin{equation}\label{eq:nnf-physical-embedding}
    \mathcal C = \mathcal T_{\mathrm{co}}\circ\mathcal L_A\circ\mathcal B,
    \qquad \vec z = \mathcal C(\vec u),\qquad
    \mathcal M\circ\mathcal C = \mathcal C\circ\mathcal N\,.
\end{equation}
The last equation is the conjugacy equation.  Tracking a normal-coordinate
point and then transforming it gives the same result as transforming first
and tracking in physical coordinates, to the retained order.

\begin{figure}[htbp]
\centering
\begin{tikzpicture}[>=stealth,node distance=1.6cm and 4.8cm]
 \node (un) {$\vec u_n$};
 \node[right=of un] (unp) {$\vec u_{n + 1}$};
 \node[below=of un] (zn) {$\vec z_n$};
 \node[below=of unp] (znp) {$\vec z_{n + 1}$};
 \draw[->,thick] (un) -- (unp)
   node[midway,above] {$\mathcal N$: normal-form dynamics};
 \draw[->,thick] (zn) -- (znp)
   node[midway,below] {$\mathcal M$: physical tracking};
 \draw[->,thick] (un) -- (zn) node[midway,left] {$\mathcal C$};
 \draw[->,thick] (unp) -- (znp) node[midway,right] {$\mathcal C$};
\end{tikzpicture}
\caption{The embedding $\mathcal C$ describes the physical geometry.
The map $\mathcal N$ describes the dynamics in normal coordinates.
Both paths give the same physical point to the retained order.}
\label{fig:nnf-conjugacy}
\end{figure}


\subsection{Factor the residual map into Lie maps}

For clarity, define the Hamiltonian point map
\begin{equation}\label{eq:nnf-hamiltonian-point-map}
    \mathcal E_H(\vec u) = \e^{-\lie{H}}\vec u,
    \qquad \widehat{\mathcal E_H} = \e^{-\lie{H}}\,.
\end{equation}
This is the unit-parameter \emph{forward} Hamiltonian flow.  The minus sign
follows from the convention $\Lie{H}g = \poisson{H}{g}$ in
Eq.~\eqref{eq:lie-flow-relation}; it is essential when extracting tune
shifts from a Hamiltonian.

Removing the linear rotation gives
$\mathcal R^{-1}\circ\mathcal F = \operatorname{id} + \mathcal O(\|\vec w\|^2)$,
where $\mathcal R(\vec w) = R\vec w$.  In the local canonical neighbourhood
used here, the formal logarithm of this near-identity symplectic map is a
Hamiltonian vector field (with its Hamiltonian defined up to an irrelevant
constant).  This is the local setting of the Dragt--Finn Lie factorisation of analytic
symplectic maps~\cite{DragtFinn1976}.  We may therefore write, order by order,
\begin{equation}\label{eq:nnf-residual-log}
    \mathcal F = \mathcal R\circ\mathcal E_K,
    \qquad K = K_3 + K_4 + K_5 + \cdots\,.
\end{equation}
Here $K_r$ is a homogeneous polynomial of degree $r$ in the canonical
coordinates.  Equivalently, a degree-ordered Lie factorisation is
\begin{align}\label{eq:nnf-lie-factorisation}
    \mathcal F
    & = \mathcal R\circ\cdots\circ\mathcal E_{f_5}
    \circ\mathcal E_{f_4}\circ\mathcal E_{f_3},\\
    \widehat{\mathcal F}
    & = \e^{-\lie{f_3}}\e^{-\lie{f_4}}
    \e^{-\lie{f_5}}\cdots\widehat{\mathcal R}\,.
\end{align}
Point-map composition and substitution-operator multiplication have the
opposite order, exactly as in Eq.~\eqref{eq:pullback-order}.

A degree-$r$ Hamiltonian first changes coordinates at degree $r-1$ because
the Poisson bracket differentiates the Hamiltonian once.  A Taylor map
known through coordinate degree $d$ therefore supplies generators through
degree $d + 1$.  Cubic generators include sextupole effects, and quartic
generators include octupole effects, but these are degrees of the
\emph{composed and centred one-turn map}, not labels of individual magnets.
For example, brackets of two cubic generators are quartic:
\begin{equation}
    \deg\poisson{f_r}{g_s} = r + s-2\,.
\end{equation}
Consequently, sextupoles also contribute to quartic detuning and resonance
terms at second order in their strengths.



\section{Solving the Normal-Form Problem Order by Order}

\subsection{The resonance basis diagonalises the homological operator}

Write $\vec u = (Q_1,P_1,\ldots,Q_n,P_n)$ and introduce the phasors
\begin{equation}\label{eq:nnf-phasors}
    a_j = \frac{Q_j-\ii P_j}{\sqrt 2} = \sqrt{J_j}\,\e^{\ii\phi_j},\qquad
    J_j = a_ja_j^* = \frac{Q_j^2 + P_j^2}{2}\,.
\end{equation}
With this phase convention,
$Q_j = \sqrt{2J_j}\cos\phi_j$,
$P_j = -\sqrt{2J_j}\sin\phi_j$ and $\poisson{\phi_j}{J_k} = \delta_{jk}$.
The linear map advances $a_j$ by $\e^{\ii\mu_j}$.

For multi-indices $\vec p,\vec q$ of nonnegative integers, use monomials
\begin{equation}
    m_{\vec p\vec q} = \prod_{j = 1}^n a_j^{p_j}(a_j^*)^{q_j} = \prod_{j = 1}^n J_j^{(p_j + q_j)/2}
    \e^{\ii(\vec p-\vec q)\cdot\vec\phi}\,.
\end{equation}
They are eigenfunctions of the substitution operator:
\begin{equation}\label{eq:nnf-resonance-eigenvalue}
    \widehat{\mathcal R}m_{\vec p\vec q} = \e^{\ii\vec k\cdot\vec\mu}m_{\vec p\vec q},
    \qquad \vec k = \vec p-\vec q\,.
\end{equation}
The same eigenvalue problem that diagonalised the linear matrix now
diagonalises an operator on polynomials.  The relevant eigenvalues of
$1-\widehat{\mathcal R}$ are
$1-\e^{\ii\vec k\cdot\vec\mu}$.  They vanish at
\begin{equation}\label{eq:nnf-resonance-condition}
    \vec k\cdot\vec\nu\in\mathbb Z,\qquad \vec\nu = \frac{\vec\mu}{2\pi}\,.
\end{equation}
This is a resonance of a \emph{discrete one-turn map}; an integer multiple
of $2\pi$ is as resonant as zero phase advance.


\subsection{Derivation of the homological equation}

At the first nonlinear order, take $\mathcal B = \mathcal E_G$ and
$\mathcal F = \mathcal R\circ\mathcal E_K$.  Pulling back
Eq.~\eqref{eq:nnf-conjugacy} gives
\begin{equation}
    \widehat{\mathcal N} = \e^{-\lie{G}}\e^{-\lie{K}}\widehat{\mathcal R}\e^{\lie{G}}\,.
\end{equation}
Since $\mathcal R$ is canonical,
\begin{equation}
    \widehat{\mathcal R}\,\Lie{G}\,
    \widehat{\mathcal R}^{-1} = \Lie{G\circ\mathcal R}\,.
\end{equation}
Combining the exponentials to leading order therefore yields
\begin{equation}\label{eq:nnf-leading-conjugation}
    \widehat{\mathcal N} = \e^{-\lie{K + G-G\circ\mathcal R + \cdots}}\widehat{\mathcal R}\,.
\end{equation}
The omitted Poisson brackets belong to higher polynomial degrees.  After
lower orders have been normalised, precisely the same calculation applies
to the residual degree-$r$ generator $K_r$.

Let $\mathcal P$ project onto the terms that will be retained: always the
action terms, and also any selected resonant or near-resonant harmonics.
Every exact-resonant term present through the retained order must be
included, since its denominator vanishes.  Optional near-resonant terms
are included together with their conjugate harmonics.
Set $H_r = \mathcal P K_r$.  The degree-$r$ homological equation is
\begin{equation}\label{eq:nnf-homological-equation}
    (1-\widehat{\mathcal R})G_r = H_r-K_r\,.
\end{equation}
If $K_r = \sum K_{\vec p\vec q}m_{\vec p\vec q}$, its solution is
\begin{equation}\label{eq:nnf-homological-solution}
    G_{\vec p\vec q} = \begin{cases}
    \displaystyle-\frac{K_{\vec p\vec q}}
    {1-\e^{\ii(\vec p-\vec q)\cdot\vec\mu}},
    &m_{\vec p\vec q}\text{ is removed},\\[1.0ex]
    0, &m_{\vec p\vec q}\text{ is retained}.
    \end{cases}\,.
\end{equation}
Setting the retained part of $G_r$ to zero fixes a normalisation convention;
adding a commuting transformation would otherwise give a different
embedding with the same normal-form dynamics.  Conjugate coefficients must
be paired so that $G_r$ remains real in real coordinates.

Applying $\mathcal E_{G_r}^{-1}$ on the left and
$\mathcal E_{G_r}$ on the right removes the chosen degree-$r$ terms.
Recompute the higher-degree residual and repeat.  The resulting
normalisation is a composition of near-identity canonical maps, rather
than a simple sum of the $G_r$.  The Baker--Campbell--Hausdorff brackets
generated during these steps are needed for the higher-order coefficients.


\subsection{What is removed, and what remains}

When $\vec p = \vec q$, the monomial is $\prod_jJ_j^{p_j}$ and its
eigenvalue is one for every tune.  These terms cannot be removed by the
homological equation.  They give the part of the normal-form Hamiltonian
that depends only on actions.  Away from other resonances, they are the
only terms retained.

Terms with $\vec p\ne\vec q$ describe angular harmonics.  Their coefficients
enter the nonlinear deformation through $G_r$ if they are nonresonant,
and remain in the normal-form dynamics if they are resonant.  For example,
$a_x^3$ has phase $3\phi_x$ and drives $3\nu_x\in\mathbb Z$, whereas
$a_x^2(a_y^*)^2$ has phase $2\phi_x-2\phi_y$ and drives
$2\nu_x-2\nu_y\in\mathbb Z$.  Calling a term nonresonant does not mean that
its physical effect vanishes: it has been transferred into the coordinate
transformation.

The normal form is not unique.  The homological equation determines only the
part of $G_r$ outside the kernel of $1-\widehat{\mathcal R}$.  Adding a
commuting generator changes the embedding $\mathcal B$ while leaving the
normal-form dynamics equivalent.  Setting the retained part of $G_r$ to zero,
as in Eq.~\eqref{eq:nnf-homological-solution}, is therefore a gauge choice for
the normalising transformation, analogous to fixing a phase convention in the
linear eigenvectors.

Near a resonance, a small denominator can make the proposed transformation
large.  Such a term should be retained in a resonant normal form.  The
decision depends on its strength and the amplitudes of interest, as well
as the distance to the resonance; a denominator threshold alone is a
numerical policy, not a universal physical criterion.



\section{Geometry of Tori and Dynamics on Tori}

\subsection{An action-only normal form}

If, through the retained polynomial order, the only harmonics in the kernel
of $1-\widehat{\mathcal R}$ are the action monomials (and no additional
near-resonant harmonics have deliberately been retained), the transformed
map is
\begin{equation}\label{eq:nnf-action-normal-form}
    \mathcal N = \mathcal R\circ\mathcal E_{h(\vec J)} = \mathcal E_{H_{\mathrm{NF}}},\qquad
    H_{\mathrm{NF}}(\vec J) = \sum_j\mu_jJ_j + h(\vec J)\,.
\end{equation}
The two factors commute because both Hamiltonians depend only on actions.
Hamilton's equations give
\begin{align}\label{eq:normal_one_turn}
    J_{j,n + 1}& = J_{j,n},\\
    \phi_{j,n + 1}& = \phi_{j,n} + \Omega_j(\vec J),\qquad
    \Omega_j = \frac{\partial H_{\mathrm{NF}}}{\partial J_j}\,.
\end{align}
These relations hold exactly for the retained action-only normal form and
up to the truncation error for the original map.  The nonlinear tune is
\begin{equation}\label{eq:nnf-tune}
    \nu_j(\vec J) = \frac{\Omega_j(\vec J)}{2\pi}\,.
\end{equation}
As in the linear case, a one-turn phase fixes the tune modulo an integer;
the branch must be followed continuously from the linear optics.  The function
$H_{\mathrm{NF}}$ is an \emph{interpolating Hamiltonian} for the one-turn
normal-form map: it generates the same retained one-turn transformation.  It
is not, in general, the local $s$-dependent Hamiltonian of the accelerator.

Fixing the positive actions gives a product of circles in normal
coordinates, with the angles $\vec\phi$ locating points on that torus.
Its image under $\mathcal C$ is a generally deformed torus in physical
phase space.
If an action vanishes, the corresponding circle collapses and the torus
has lower dimension.  Approximate invariants in physical coordinates are
\begin{equation}
    I_j(\vec z) = J_j\bigl(\mathcal C^{-1}(\vec z)\bigr)\,.
\end{equation}
Computing an action directly from the \emph{linearly} normalised
coordinates misses $\mathcal B^{-1}$ and generally produces turn-dependent
oscillations rather than an invariant.

The map $\mathcal C$ specifies the shape, local size, mean displacement,
and harmonic content of the oscillation.  The Hamiltonian
$H_{\mathrm{NF}}$ specifies how its phases advance.  A term in the embedding
can substantially distort an orbit without directly changing its tune at
the same perturbative order.


\subsection{Resonant dynamics}

With resonant terms retained, write instead
\begin{equation}\label{eq:nnf-resonant-map}
    \mathcal N = \mathcal R\circ\mathcal E_{Z(\vec J,\vec\phi)},\qquad
    Z = h(\vec J) + \sum_{\vec k\ne0}
    V_{\vec k}(\vec J)\e^{\ii\vec k\cdot\vec\phi}\,,
\end{equation}
with conjugate pairs making $Z$ real.  During the residual Hamiltonian
flow, $\dot J_j = -\partial Z/\partial\phi_j$.  Individual actions can
therefore vary, and the dynamics can develop resonance islands,
separatrices, and mode exchange.  For one retained resonance vector
$\vec k$, combinations $\vec l\cdot\vec J$ with
$\vec l\cdot\vec k = 0$ are conserved by this ideal normal form, since the
angular dependence is only through $\vec k\cdot\vec\phi$.

At exact resonance the retained harmonics commute with the discrete
rotation.  Near resonance they generally do not.  One may introduce a
slow resonant angle and a rotating frame, retaining the tune mismatch
alongside the resonant Hamiltonian.  Simply replacing the map by
$\mathcal E_{\sum\mu_jJ_j + Z}$ without the required commutation or
Baker--Campbell--Hausdorff corrections is not justified.

Thus the interpretation in terms of deformed constant-action tori applies
to the nonresonant, action-only approximation.  A resonant normal form
still separates coordinates from dynamics, but its dynamics is richer
than rotation on those tori.  Section~\ref{sec:nnf-rdts} returns to the
complex coefficients commonly used in accelerator physics to quantify these
resonance harmonics.



\section{Amplitude Detuning}

Expand the nonresonant Hamiltonian as
\begin{equation}\label{eq:nnf-detuning-hamiltonian}
    H_{\mathrm{NF}} = \sum_j\mu_jJ_j + \frac12\sum_{jk}D_{jk}J_jJ_k + \frac16\sum_{jkl}E_{jkl}J_jJ_kJ_l + \cdots\,,
\end{equation}
where $D$ is symmetric and $E$ is symmetric in its indices.  Then
\begin{equation}\label{eq:nnf-detuning-tunes}
    \nu_j(\vec J) = \nu_{j0} + \frac{1}{2\pi}\sum_kD_{jk}J_k + \frac{1}{4\pi}\sum_{kl}E_{jkl}J_kJ_l + \cdots\,.
\end{equation}
The tune footprint is the image of an action domain under this frequency
map.  Its slope is the amplitude-detuning matrix.  Since the frequencies
derive from one scalar Hamiltonian,
\begin{equation}\label{eq:nnf-detuning-symmetry}
    \frac{\partial\nu_j}{\partial J_k} = \frac{\partial\nu_k}{\partial J_j}\,.
\end{equation}
This relation is a useful consistency check when the actions use the same
canonical normalisation.

Quartic action terms produce detuning linear in action, hence quadratic
in oscillation amplitude.  Sextic terms produce quadratic-in-action
detuning.  The coefficient is not obtained from the raw quartic kick
alone if lower-degree nonlinearities are present: normalising the cubic
terms also generates quartic contributions.  Similarly, the physical
oscillation amplitude is found through $\mathcal C$, so it should not be
identified with $\sqrt{2J_j}$ in physical units.



\section{Beta Functions Depending on Amplitude}

\subsection{A definition from the fundamental harmonic}

The linear beta function measures how a normalised circle projects into a
physical coordinate.  At finite amplitude, a deformed torus need not
project as an ellipse and contains higher harmonics.  There is therefore
no single nonlinear beta function that automatically has every property
of the linear Courant--Snyder beta.  A definition must state which feature
of the motion it measures.

For regular nonresonant motion, write the physical embedding at a location
$s$ as a Fourier series in the normal-form angles:
\begin{equation}\label{eq:nnf-torus-fourier}
    x(s,\vec J,\vec\phi)-x_{\mathrm{co}}(s) = \sum_{\vec k\in\mathbb Z^n}
    X_{\vec k}(s,\vec J)\e^{\ii\vec k\cdot\vec\phi},
    \qquad X_{-\vec k} = X_{\vec k}^*\,.
\end{equation}
Let $\vec e_j$ be the unit harmonic vector of mode $j$.  A beta function
based on its fundamental position harmonic is
\begin{equation}\label{eq:nnf-amplitude-beta}
    \beta_{x,j}^{(1)}(s,\vec J) = \frac{2|X_{\vec e_j}(s,\vec J)|^2}{J_j}\,.
\end{equation}
The zero-action value is understood as a limit.  For uncoupled linear
horizontal motion, $X_{\vec e_x} = \sqrt{\beta_xJ_x/2}\,
\e^{\ii\psi_x(s)}$, so Eq.~\eqref{eq:nnf-amplitude-beta} reduces to the
usual $\beta_x(s)$.  For coupled linear motion, it reduces to the projection
beta of that normal mode into $x$.

The nonlinear correction has the form
\begin{equation}
    \beta_{x,j}^{(1)}(s,\vec J) = \beta_{x,j}^{\mathrm{lin}}(s) + \sum_k b_{x,jk}(s)J_k + \cdots\,.
\end{equation}
Its coefficients come from the normalising map $\mathcal B$ and the
transport of the embedding to $s$, whereas the detuning coefficients come
from $H_{\mathrm{NF}}$.  These are different quantities.  In particular,
the linear relation between beta and local phase advance does not become
a general nonlinear identity by substituting an amplitude-dependent beta.


\subsection{Mean displacements, higher harmonics, and transport}

The coefficient $X_0$ can be nonzero: the mean position of a finite-action
torus can differ from the zero-action closed orbit.  It is part of the
nonlinear embedding and should not be mistaken for an error in the
closed-orbit subtraction.  Higher $X_{\vec k}$ describe additional
distortion.  In an uncoupled one-mode example, an alternative measure is
\begin{equation}
    \beta_x^{\mathrm{rms}}(J) = \frac{\left\langle(x-\langle x\rangle)^2\right\rangle_\phi}{J} = \frac{1}{J}\sum_{k\ne0}|X_k(J)|^2\,.
\end{equation}
It agrees with the linear beta, but includes all harmonics and generally
differs from $\beta_x^{(1)}$ at finite amplitude.

If $\mathcal P_{s\leftarrow s_0}$ transports physical points from the
normal-form observation point $s_0$ to $s$, the torus embedding there is
\begin{equation}
    \mathcal C_s = \mathcal P_{s\leftarrow s_0}\circ\mathcal C_{s_0}\,.
\end{equation}
A local phase convention can subsequently be chosen for the normal
coordinates.  Extracting Fourier coefficients from this transported map
gives amplitude-dependent optics along the lattice.  Fitting a single
ellipse to a strongly distorted projection does not in general reproduce
these harmonic beta functions.



\section{Resonance Driving Terms}\label{sec:nnf-rdts}

The resonance discussion above tells us \emph{which} angular harmonics are
associated with a resonance.  In accelerator physics one usually also wants a
complex number that quantifies how strongly a particular harmonic is excited
and with which phase.  These coefficients are called \emph{resonance driving
terms} (RDTs).  They are widely used to diagnose nonlinear lattices, compare
magnet configurations, and connect normal-form calculations with spectral
lines measured from turn-by-turn data~\cite{PhysRevSTAB.17.074001}.

There is a small terminology warning.  The name ``resonance driving term'' is
not completely uniform in the literature.  Some authors use it for a
coefficient of the Hamiltonian itself, while others use it for the coefficient
of the generator of the normalising transformation.  In this chapter we use
\emph{RDT} for the latter, because it follows directly from the normal-form
transformation $\mathcal B$.  Coefficients that remain in the resonant
Hamiltonian $Z$ will instead be called \emph{resonant Hamiltonian
coefficients}.  Keeping this distinction explicit avoids a common source of
factor and sign confusion.


\subsection{RDTs from the normalising transformation}

The order-by-order construction above writes $\mathcal B$ as a product of
near-identity Lie maps.  Through any fixed truncation order, these maps may be
combined with the Baker--Campbell--Hausdorff series into a single formal
near-identity map
\begin{equation}\label{eq:nnf-rdt-generator}
    \mathcal B = \mathcal E_F = \e^{-:F:},
    \qquad
    F = \sum_{\vec p,\vec q} f_{\vec p\vec q}
    m_{\vec p\vec q}\,.
\end{equation}
The complex coefficients $f_{\vec p\vec q}$ are the RDTs in the convention
used here.  At the first order at which a monomial appears, $F$ and the
corresponding order-by-order generator $G_r$ coincide.  Equation
\eqref{eq:nnf-homological-solution} therefore gives immediately, for a term
that is removed from the normal-form dynamics,
\begin{equation}\label{eq:nnf-rdt-from-k}
    f_{\vec p\vec q}^{(r)}
    = G_{\vec p\vec q}
    = -\frac{K_{\vec p\vec q}}
    {1-\e^{\ii\vec k\cdot\vec\mu}},
    \qquad
    \vec k = \vec p-\vec q\,.
\end{equation}
At higher perturbative orders, $F$ also contains the BCH combinations of the
lower-order normalising generators, so the full RDT is not in general just a
raw coefficient of one magnet or of one Taylor-map order.

For two transverse degrees of freedom it is customary to write
\begin{equation}\label{eq:nnf-rdt-jklm}
    F = \sum_{jklm} f_{jklm}
    a_x^j(a_x^*)^k a_y^l(a_y^*)^m\,.
\end{equation}
The harmonic vector of this monomial is
\begin{equation}
    \vec k = (j-k,\,l-m)\,,
\end{equation}
so the associated resonance condition is
\begin{equation}\label{eq:nnf-rdt-resonance}
    (j-k)\nu_x + (l-m)\nu_y = N,
    \qquad N\in\mathbb Z\,.
\end{equation}
For example, $f_{3000}$ is associated with $3\nu_x=N$, while $f_{2010}$ is
associated with $2\nu_x+\nu_y=N$.  The polynomial degree
$j+k+l+m$ and the resonance order $|j-k|+|l-m|$ are related but are not the
same quantity.  Higher-degree monomials can carry the same harmonic vector and
therefore contribute to the same resonance with an amplitude-dependent
strength.

Equation~\eqref{eq:nnf-rdt-from-k} also shows why an RDT becomes large close
to its resonance.  Since
\begin{equation}
    1-\e^{\ii\theta}
    = -2\ii\e^{\ii\theta/2}\sin\!\left(\frac{\theta}{2}\right)\,,
\end{equation}
the normalising transformation contains a small denominator when
$\vec k\cdot\vec\mu$ approaches an integer multiple of $2\pi$.  The RDT
therefore measures not only the strength of the corresponding harmonic in the
one-turn perturbation, but also how strongly repeated turns coherently build
up that harmonic at the chosen tunes.


\subsection{Physical meaning of the complex coefficient}

The modulus and phase of an RDT have complementary meanings.  Roughly
speaking, $|f_{\vec p\vec q}|$ measures the size of the corresponding
nonlinear distortion, while its complex phase fixes the phase of that
harmonic relative to the chosen observation point and phase convention.  For
a distributed lattice, the numerator $K_{\vec p\vec q}$ is, to first order in
the nonlinear strengths, a coherent sum of contributions from the nonlinear
elements.  Schematically,
\begin{equation}\label{eq:nnf-rdt-phasor-sum}
    K_{\vec p\vec q}(s_0)
    \sim \sum_m c_{\vec p\vec q,m}
    \e^{\ii\vec k\cdot\Delta\vec\phi(s_m,s_0)}\,.
\end{equation}
The coefficients $c_{\vec p\vec q,m}$ contain the local multipole strength and
optical factors.  The exact prefactors depend on the multipole and coordinate
conventions.  Equation~\eqref{eq:nnf-rdt-phasor-sum} is nevertheless the
important physical point: contributions from different magnets add as
phasors, so they can reinforce or cancel according to their phase advances.
This is why the location of nonlinear correctors matters even when their
integrated strengths are identical.

The same coefficients appear directly in the shape of the motion.  From
Eq.~\eqref{eq:nnf-rdt-generator}, a normal-form phasor $a_j$ is mapped back to
the linearly normalised physical coordinates as
\begin{equation}\label{eq:nnf-rdt-spectral-lines}
    a_j^{\mathrm{lin}}
    = \e^{-:F:}a_j^{\mathrm{NF}}
    = a_j^{\mathrm{NF}}
    + \ii\frac{\partial F}{\partial a_j^*}
    + \mathcal O(F^2)\,.
\end{equation}
Each monomial of $F$ therefore produces a definite angular harmonic in the
turn-by-turn motion.  This is the basis for extracting RDT information from
the complex Fourier spectrum of BPM data.  It also makes clear that RDTs
belong primarily to the nonlinear \emph{coordinate transformation}: they
describe how simple normal-form motion is deformed when expressed in the
linearly normalised or physical coordinates.

Several monomials can correspond to the same resonance vector $\vec k$.  At a
specified action, their combined effect is better represented by grouping them
as
\begin{equation}\label{eq:nnf-rdt-action-dependent}
    F(\vec J,\vec\phi)
    = \sum_{\vec k}F_{\vec k}(\vec J)
    \e^{\ii\vec k\cdot\vec\phi},
    \qquad
    F_{\vec k}(\vec J)
    = \sum_{\vec p-\vec q=\vec k}
    f_{\vec p\vec q}
    \prod_j J_j^{(p_j+q_j)/2}\,.
\end{equation}
Thus a single low-order RDT is a useful diagnostic, but it is not by itself a
complete finite-amplitude measure of a resonance when higher-order terms are
important.


\subsection{What happens on the resonance}

Equation~\eqref{eq:nnf-rdt-from-k} must not be interpreted literally at an
exact resonance.  When
$\vec k\cdot\vec\mu = 2\pi N$, the denominator vanishes because the
corresponding harmonic cannot be removed by a near-identity transformation.
The divergence is a failure of the \emph{nonresonant coordinate change}, not
an infinite physical force.

In that situation the harmonic is retained in the resonant normal form,
\begin{equation}
    Z(\vec J,\vec\phi)
    \supset V_{\vec k}(\vec J)
    \e^{\ii\vec k\cdot\vec\phi} + \text{c.c.}\,,
\end{equation}
and the finite coefficient $V_{\vec k}$ is the quantity that directly enters
the resonant Hamiltonian dynamics.  Together with the distance from resonance
and the action-dependent detuning, it determines features such as island
widths, fixed points, separatrices, and action exchange.  Close to a
resonance, therefore, a large nonresonant RDT is best understood as a warning
that the corresponding harmonic should be promoted from the coordinate
transformation into a resonant normal form.

This distinction is useful in practice: away from resonance, RDTs provide a
compact description of nonlinear distortion and measurable spectral lines;
near or on resonance, the retained resonant Hamiltonian gives the more direct
description of the dynamics.



\section{Momentum Dependence and Nonlinear Chromaticity}

\subsection{A family of off-momentum normal forms}

Suppose the relative momentum offset $\delta_0$ is constant during the
motion, as in a transverse coasting-beam model.  It then labels a family
of four-dimensional symplectic maps,
\begin{equation}
    \mathcal M_{\delta_0},\quad
    \vec z_{\mathrm{co}}(\delta_0),\quad A(\delta_0),\quad
    \mathcal B_{\delta_0},\quad H_{\mathrm{NF}}(\vec J;\delta_0)\,.
\end{equation}
Both the coordinates and the dynamics depend on momentum.  One first
solves for the off-momentum closed orbit, including its nonlinear
dispersion, and then normalises about that orbit.  Subtracting only the
on-momentum orbit, or only its linear dispersion correction, is
insufficient for a higher-order chromatic calculation.

Use the same canonical coordinate and momentum convention throughout the
family.  In particular, comparisons at fixed action require a consistent
normalisation; holding a physical launch displacement fixed generally
changes the action when $A$ or $\mathcal B$ changes.

The nonresonant Hamiltonian becomes
\begin{equation}\label{eq:nnf-chromatic-hamiltonian}
    H_{\mathrm{NF}}(\vec J;\delta_0) = \sum_j2\pi \nu_{j0}(\delta_0)J_j + \frac12\sum_{jk}D_{jk}(\delta_0)J_jJ_k + \cdots\,.
\end{equation}
Consequently,
\begin{equation}\label{eq:nnf-amplitude-chromatic-tune}
    \nu_j(\vec J;\delta_0) = \nu_{j0}(\delta_0) + \frac{1}{2\pi}\sum_kD_{jk}(\delta_0)J_k + \cdots\,.
\end{equation}
The Fourier beta functions also depend on $\delta_0$, through
$\mathcal C_s(\vec J,\vec\phi;\delta_0)$.


\subsection{Chromatic derivatives and mixed terms}

Define the chromaticity of order $r$ by derivatives, so the factorial
convention is explicit:
\begin{align}\label{eq:nnf-chromaticity-definition}
    \xi_j^{(r)}(\vec J)
    & = \left.\frac{\partial^r \nu_j(\vec J;\delta_0)}
    {\partial\delta_0^r}\right|_{\delta_0 = 0},\\
    \nu_j(\vec J;\delta_0)
    & = \nu_j(\vec J;0) + \sum_{r\ge1}
    \frac{\xi_j^{(r)}(\vec J)}{r!}\delta_0^r\,.
\end{align}
At zero action, $\xi_j^{(1)}(0)$ is the usual linear chromaticity.
The derivatives with $r\ge2$ describe nonlinear chromaticity of the
small-amplitude tune.  At finite action, even the first chromatic
derivative can depend on amplitude:
\begin{equation}
    \xi_j^{(1)}(\vec J) = \xi_j^{(1)}(0) + \frac{1}{2\pi}\sum_kD_{jk}'(0)J_k + \cdots\,.
\end{equation}
Thus nonlinear chromaticity and amplitude-dependent chromaticity are
related but distinct effects.  Mixed terms such as $J_jJ_k\delta_0$ and
$J_j\delta_0^2$ in the Hamiltonian retain both kinds of information.

If RF motion makes momentum dynamical, $\delta$ cannot be treated as a
constant external label.  The appropriate construction uses all three
canonical pairs about a six-dimensional periodic orbit.  Its nonresonant
frequencies depend on the longitudinal action $J_s$, and its resonance
basis includes $k_x\nu_x + k_y\nu_y + k_s\nu_s\in\mathbb Z$.  A scan in constant
$\delta_0$ is useful transverse optics information, but it does not itself
describe synchrobetatron modulation or synchrobetatron resonances.  The
longitudinal coordinates must also be a canonical pair in the adopted
tracking convention.



\section{A Magnet Strength as an External Parameter}

\subsection{Parameter-dependent coordinates and frequencies}

Let $\lambda$ be a magnet strength or another lattice knob.  The entire
construction becomes a family:
\begin{equation}\label{eq:nnf-parameter-family}
    \mathcal M_{\delta_0,\lambda} = \mathcal C_{\delta_0,\lambda}\circ
    \mathcal N_{\delta_0,\lambda}\circ
    \mathcal C_{\delta_0,\lambda}^{-1}\,.
\end{equation}
For example, with $\eta = \lambda-\lambda_*$ one may expand
\begin{equation}
    D_{jk}(\delta_0,\lambda) = D_{jk,*} + D_{jk,\delta}\delta_0 + D_{jk,\lambda}\eta + D_{jk,\delta\lambda}\delta_0\eta + \cdots\,.
\end{equation}
The embedding has its own expansion.  A knob can therefore change both
the tune footprint and the amplitude-dependent optics.

External parameters are not canonical dynamical variables: they have no
conjugate momentum in this calculation and do not enter the transverse
Poisson bracket.  Taylor coefficients in these parameters can be carried
through the map operations as well.  The normal-form sequence follows
polynomial degree in the dynamical coordinates; parameter degree can be
truncated separately.  The resonance denominators, eigenvectors, and
closed orbit must also be expanded or differentiated, rather than held
fixed without justification.

Choose mode labels, eigenvector phases, and tune branches continuously
along the parameter family.  If a selected denominator crosses a
resonance, a nonresonant normalisation cannot be differentiated smoothly
through that crossing; retain the relevant harmonic or use a different
local expansion.


\subsection{Example: an octupole at a specified location}

Consider a thin normal octupole at $s_o$, with integrated strength
$\kappa = \kappa_* + \lambda$ defined by the Hamiltonian
\begin{equation}\label{eq:nnf-octupole-hamiltonian}
    H_o(x,y) = \frac{\kappa}{24}(x^4-6x^2y^2 + y^4)\,.
\end{equation}
Its point kick leaves $x,y$ fixed and gives
\begin{equation}
    \Delta p_x = -\frac{\kappa}{6}(x^3-3xy^2),\qquad
    \Delta p_y = \frac{\kappa}{6}(3x^2y-y^3)\,.
\end{equation}
The Hamiltonian and the strength use the same canonical momentum
convention as the lattice map.  This explicit definition avoids
factorial and sign ambiguities in multipole conventions.

First assume an uncoupled linear lattice, a closed orbit through the
octupole centre, no other nonlinearities, and weak $\kappa$, away from
the relevant resonances.  At the octupole,
\begin{equation}
    x = \sqrt{2\beta_x(s_o)J_x}\cos\theta_x,\qquad
    y = \sqrt{2\beta_y(s_o)J_y}\cos\theta_y\,.
\end{equation}
The local angles include the phase advance from the observation point.
The quartic action part is the angle average of $H_o$:
\begin{align}\label{eq:nnf-octupole-average}
    \langle x^4\rangle& = \frac32\beta_x^2J_x^2,
    &\langle x^2y^2\rangle& = \beta_x\beta_yJ_xJ_y,\\
    h_o& = \kappa\left(
    \frac{\beta_x^2J_x^2}{16}
    -\frac{\beta_x\beta_yJ_xJ_y}{4} + \frac{\beta_y^2J_y^2}{16}\right)\,.
\end{align}
Here the beta functions are evaluated at $s_o$.  Differentiating gives
the first-order tune shifts
\begin{align}\label{eq:nnf-octupole-detuning}
    \Delta\nu_x& = \frac{\kappa}{16\pi}\beta_x^2J_x
    -\frac{\kappa}{8\pi}\beta_x\beta_yJ_y,\\
    \Delta\nu_y& = \frac{\kappa}{16\pi}\beta_y^2J_y
    -\frac{\kappa}{8\pi}\beta_x\beta_yJ_x\,.
\end{align}
For this example, $\partial\Delta\nu_j/\partial\lambda$ is obtained by
replacing $\kappa$ with one.  An on-axis octupole changes the finite-action
tunes at this order while leaving the zero-action linear tunes unchanged.

The same octupole also deforms the tori.  To see both outputs explicitly,
set $J_y = 0$, take the observation point immediately before the kick, and
write the remainder of the linear turn as $\mathcal R$.  Then
$\mathcal F = \mathcal R\circ\mathcal E_{K_4}$ with
\begin{equation}\label{eq:nnf-octupole-harmonics}
    K_4 = \kappa\beta_x^2J_x^2
    \left(\frac{1}{16} + \frac{\cos2\phi_x}{12} + \frac{\cos4\phi_x}{48}\right)\,.
\end{equation}
The constant harmonic is $h_o$.  Away from the second- and fourth-order
resonances, the other two harmonics give
\begin{equation}\label{eq:nnf-octupole-normaliser}
    G_4 = \frac{\kappa\beta_x^2J_x^2}{24\sin\mu_x}
    \sin(2\phi_x-\mu_x) + \frac{\kappa\beta_x^2J_x^2}{96\sin2\mu_x}
    \sin(4\phi_x-2\mu_x)\,.
\end{equation}
Each harmonic solves
$G_4(\phi_x)-G_4(\phi_x + \mu_x) = h_o-K_4$.
The map $\mathcal B = \mathcal E_{G_4}$ gives
\begin{equation}
    \vec w = \vec u-\poisson{G_4}{\vec u} + \mathcal O(\kappa^2)\,,
\end{equation}
at this perturbative order.  The fundamental position harmonic yields,
at this particular observation point and with this phase convention,
\begin{equation}\label{eq:nnf-octupole-beta}
    \beta_x^{(1)}(J_x) = \beta_x-\frac{\kappa\beta_x^3}{8}\cot\mu_x\,J_x + \mathcal O(\kappa^2)\,,
\end{equation}
at fixed sufficiently small $J_x$.  This comes from the $2\phi_x$ term
in $G_4$; other harmonics affect the orbit shape as well.  Transporting the
embedding gives the beta correction at other locations.  The detuning
and this beta correction therefore arise from different parts of the
\emph{same} normal-form calculation.

In two planes the nonconstant harmonics include $4\phi_x$,
$4\phi_y$, $2\phi_x$, $2\phi_y$, and
$2\phi_x\mathbin{\pm}2\phi_y$.  Their coefficients acquire phase
factors when transported to a different observation point.  Contributions
from several octupoles add as phasors for a specified resonance, so their
relative locations matter.  Their first-order action averages instead add
without these phase factors, weighted by the local beta functions.

An off-axis orbit produces feed-down.  For example, take
$x_{\mathrm{co}}(s_o,\delta_0) = D_x(s_o)\delta_0 + \mathcal O(\delta_0^2)$
and $y_{\mathrm{co}} = 0$.  The quadratic Hamiltonian about that orbit
contains the effective integrated quadrupole strength
\begin{equation}
    \Delta k_1 = \frac{\kappa}{2}D_x^2\delta_0^2 + \mathcal O(\delta_0^3)\,.
\end{equation}
To first order in $\kappa$, its small-amplitude tune contribution is
\begin{equation}\label{eq:nnf-octupole-chromaticity}
    \Delta\nu_{x0} = \frac{\kappa\beta_xD_x^2}{8\pi}\delta_0^2 + \mathcal O(\delta_0^3),\qquad
    \Delta\nu_{y0} = -\frac{\kappa\beta_yD_x^2}{8\pi}\delta_0^2 + \mathcal O(\delta_0^3)\,.
\end{equation}
Here the beta and dispersion are their on-momentum values at $s_o$.
Thus this feed-down gives a second chromatic derivative even though the
on-axis, on-momentum linear tune response vanishes.  Momentum dependence
of the beta functions and of the strength convention also changes the
detuning coefficients in Eq.~\eqref{eq:nnf-octupole-detuning}, producing
mixed amplitude--chromatic effects.  In a lattice with existing sextupoles,
offsets, or coupling, the full centred map includes additional terms and
must be normalised rather than using these isolated-octupole formulae
unchanged.



\section{Example Implementation in Pseudocode}

The algorithm below assumes access to Taylor maps, composition, inversion,
and Poisson brackets.  A TPSA implementation is one way to supply these
operations; the following chapter describes that representation.  Every
map operation is truncated to coordinate degree $d$, and Hamiltonians are
kept through degree $d + 1$.  The resonance selection is fixed for the local
parameter domain.  All compositions below are compositions of
\emph{point maps}, with the rightmost map acting first.
The Taylor representation of $\mathcal M$ takes absolute physical
coordinates as arguments, even though it is expanded about the closed
orbit.  The routine \texttt{symplectic\_eigenbasis} performs the
eigenvector part of the preceding chapter's matrix algorithm.

\begin{lstlisting}[basicstyle=\ttfamily\small,columns=fullflexible,
                   keepspaces=true,frame=single]
input: one_turn_map, coordinate_degree d, parameter_orders
       fixed local set of retained resonance harmonics

co = solve_closed_orbit(one_turn_map, delta0, lambda)
M  = taylor_map(one_turn_map, around=co,
                coordinate_degree=d, parameter_orders=parameter_orders)
A, mu = symplectic_eigenbasis(jacobian(M, co))
T = translation(co)
L = linear_map(A)
F = inverse(L) o inverse(T) o M o T o L
R = rotation_map(mu)

current = F
B = identity_map()
Z = 0                       # retained nonlinear Hamiltonian

for r = 3, ..., d + 1:
    # All lower degrees have already been normalised.
    residual = inverse(E(Z)) o inverse(R) o current
    V = homogeneous_part(residual - identity_map(), r - 1)
    K_r = -dot(u, S * V) / r # V = S * gradient(K_r)
    check_gradient(K_r, -S * V)

    Kc = to_phasor_basis(K_r)
    Hc = 0
    Gc = 0
    for coefficient c, monomial a**p * conjugate(a)**q in Kc:
        k = p - q
        if k == 0 or harmonic_is_retained(k):
            Hc += c * monomial
        else:
            denom = 1 - exp(i * dot(k, mu))
            check_small_denominator(c, denom, action_domain)
            Gc += (-c / denom) * monomial

    H_r = to_real_canonical_coordinates(Hc)
    G_r = to_real_canonical_coordinates(Gc)
    check_reality(H_r, G_r)
    step = E(G_r)            # E(H)(u) = exp(-:H:) u
    current = inverse(step) o current o step
    B = B o step
    Z = Z + H_r
    # Recomputed residuals include all higher-order brackets.
\end{lstlisting}

The generator extraction uses Euler's identity for a homogeneous
polynomial: if $\vec V_{r-1} = S\nabla K_r$, then
$K_r = -\vec u^TS\vec V_{r-1}/r$.  The gradient check detects a residual
that is not Hamiltonian to the required order.  In the residual above,
$\mathcal E(Z)^{-1}$ removes the already retained lower-order flow;
brackets with a new degree-$r$ term affect higher degrees and are picked
up on subsequent iterations.  All exponentials, including that of $Z$,
must be evaluated with the same truncation rules.

If parameters are represented by Taylor series, the closed-orbit solve,
linear normalisation, and resonance denominators are performed with those
series as coefficients.  Polynomial degrees in the loop refer only to
the dynamical coordinates.  Alternatively, repeat the calculation at
nearby parameter values while keeping phase gauges and mode labels
consistent.  The latter approach needs numerical convergence checks for
the resulting derivatives.

The following steps check both the conjugacy and the form of the retained
dynamics.  Invariant-torus optics and tune derivatives are then extracted
when the retained Hamiltonian depends only on actions.

\begin{minipage}{\linewidth}
\begin{lstlisting}[basicstyle=\ttfamily\small,columns=fullflexible,
                   keepspaces=true,frame=single]
N = current
C = T o L o B
check_coefficients(M o C - C o N, through_degree=d)
check_coefficients(N - R o E(Z), through_degree=d)
check_symplecticity(C, N, through_available_order=d - 1)

if only_action_terms_are_retained:
    H_NF = dot(mu, J) + Z
    tunes = gradient(H_NF, J) / (2*pi)
    chromaticities = derivatives(tunes, delta0, at=0)
    knob_responses = derivatives(tunes, lambda, at=lambda_star)
    for observation_location s:
        C_s = transport_map(s, from=s0) o C
        X = fourier_coefficients_of_embedding(C_s)
        beta_x_mode_j = 2 * abs(X[unit_harmonic(j)])**2 / J[j]

return C, N, Z, observables_if_nonresonant
\end{lstlisting}
\end{minipage}



\section{Validity and Practical Checks}

A finite-order normal form describes a neighbourhood of the chosen
closed orbit.  Its useful amplitude range depends on the nonlinear
strengths, omitted orders, and resonance denominators.  Formal existence
of the series does not guarantee convergence.  Under the additional
nondegeneracy and nonresonance assumptions of KAM theory, many sufficiently
small invariant tori survive a small Hamiltonian perturbation, but this
does not imply that every action in a finite domain labels an exact torus.

Several checks directly test what the calculation claims.  The
conjugacy residual $\mathcal M\circ\mathcal C-\mathcal C\circ\mathcal N$
should have no coefficients through the retained coordinate degree.
The computed maps should satisfy the symplectic Jacobian condition
through the order available after differentiating the Taylor map.
For an action-only normal form, track particles launched through
$\mathcal C$ and compare measured tunes, harmonic amplitudes, and the
variation of $I_j$ with the predicted quantities.  Increasing the map
order or decreasing the actions should reduce discrepancies where the
expansion is useful.

Near a separatrix or in a chaotic region, a single quasiperiodic torus
embedding and its amplitude-dependent beta are no longer an adequate
description.  Dynamic aperture and long-term stability therefore require
tracking and other diagnostics in addition to local normal forms.
The practical value of the normal form is that it explains the local
frequency shifts, orbit deformations, and resonances in one consistent
canonical construction.
